% English translation of checked_pages/esq_119_021.tex; source: cote 119, site 21, PDF page 22, typescript page 19.
% AI-translated by GPT-6 Astra (OpenAI, ChatGPT Pro mode).
\scanpage{119/21}
\mspage{19}

of the euclidean or hyperbolic plane (*). The connection
of the combinatorial theory with the ``conformal'' theory of regular
tessellations of the hyperbolic plane was sensed before
the connection of finite maps with finite coverings of the
projective line appeared. Once this connection is understood, it becomes obvious
that it must also extend to infinite maps (regular
or otherwise): \uline{every map, finite or otherwise, is canonically realized on
a conformal surface} (compact if and only if the map is
finite), \uline{as a ramified covering of the complex projective line,
ramified only at the points} 0, 1, $\infty$. The only difficulty
here was to work out the dictionary between topological
maps and sets with operators, which posed some conceptual
problems in the infinite case, beginning with the very notion of
a ``topological map''. In particular, both for the
internal coherence of the dictionary and so as not to miss
certain interesting cases of infinite maps, it proves necessary
not to exclude vertices and faces of infinite order. This foundational
work was also carried out by J. Malgoire and C. Voisin, following
on from their first work on finite maps, and their theory
indeed supplies everything one was entitled to expect from it
(and even more...).

It was in 1977 and 1978, alongside two C 4 courses\tnote{TN-021-01}{``C 4'' is the course designation used in the source; it is retained here.}
on the geometry of the cube and that of the icosahedron, that I
began to take an interest in regular polyhedra, which then appeared
to me as particularly concrete ``geometric realizations''
of combinatorial maps, the vertices, edges and faces being
realized respectively as points, lines and planes
in a suitable three-dimensional affine space, with the incidence
relations respected. This notion of geometric realization

\rule{4cm}{0.4pt}

(*) In these statements, one should not exclude the case in which $p$, $q$
can take the value $+\infty$, which one encounters in particular very
naturally in tessellations associated with certain infinite regular
polyhedra, cf. below.

\margin{21}

